
\documentclass[11pt,a4paper]{article}
\usepackage[spanish]{babel}
\usepackage[utf8]{inputenc}
\usepackage[T1]{fontenc}
\usepackage{amsmath,amssymb}
\usepackage{siunitx}
\usepackage{geometry}
\geometry{margin=2.5cm}

\begin{document}

\title{Criterios para la elección del modelo ideal en flujo compresible en conductos}
\author{}
\date{}
\maketitle

\section*{1.\ Punto de partida: expresión general para \texorpdfstring{$w/A$}{w/A}}

Del balance de energía mecánica en una conducción horizontal de sección constante, se obtiene para flujo estacionario compresible:

\begin{equation}
\left(\frac{w}{A}\right)^2 \ln\!\frac{v_2}{v_1}
\;+\;\int_1^2 \frac{\mathrm dP}{v}
\;+\;2 f \frac{L}{D}\left(\frac{w}{A}\right)^2 = 0
\end{equation}

de donde

\begin{equation}
\frac{w}{A}
=
\left[
-\frac{\displaystyle \int_1^2 \frac{\mathrm dP}{v}}
{\displaystyle \ln\!\frac{v_2}{v_1} + 2 f\frac{L}{D}}
\right]^{1/2}.
\end{equation}

Comparar modelos significa comparar el numerador
\(-\displaystyle\int_1^2 \frac{\mathrm dP}{v}\)
y el logaritmo \(\ln(v_2/v_1)\) entre distintos supuestos
(isotérmico, adiabático, con calor, etc.) manteniendo el mismo
\(\Delta P\) y el mismo \(2fL/D\).

\bigskip
\hrule
\bigskip

\section*{2.\ Isotérmico vs.\ adiabático (mismo \(P_1,P_2\))}

Suponiendo gas ideal \(Pv/T=\text{cte}\):

\begin{itemize}
  \item Isotérmico: \(T_2 = T_1\).
  \item Adiabático con caída de presión: \(T_2 < T_1\).
\end{itemize}

Para el mismo par \((P_1,P_2)\):

\begin{itemize}
  \item Se tiene \(T_{2,i} > T_{2,a}\) y por tanto \(v_{2,i} > v_{2,a}\).
  \item Entonces:
  \[
  -\int_1^2 \frac{\mathrm dP}{v}\bigg|_i
  <
  -\int_1^2 \frac{\mathrm dP}{v}\bigg|_a,
  \qquad
  \ln\frac{v_{2,i}}{v_1}
  >
  \ln\frac{v_{2,a}}{v_1}.
  \]
\end{itemize}

Con el mismo \(2fL/D\) resulta:

\begin{equation}
\boxed{
\left(\frac{w}{A}\right)_{\text{isotérmico}}
<
\left(\frac{w}{A}\right)_{\text{adiabático}}
}
\end{equation}

Es decir, el modelo adiabático da un flujo másico mayor que el isotérmico para la misma caída de presión.

\bigskip
\hrule
\bigskip

\section*{3.\ Efecto del calor respecto al caso adiabático}

Comparando casos con calor distribuido (\(Q\neq 0\)) respecto al adiabático
(\(Q=0\)), siempre con gas ideal y mismo \((P_1,P_2)\):

\begin{itemize}
  \item Si \(Q>0\) (se calienta el gas a lo largo de la tubería):
  \begin{itemize}
    \item El gas tiene \(T\) más alta y \(v\) más grande que en el adiabático.
    \item \(\Rightarrow -\displaystyle\int \frac{\mathrm dP}{v}\big|_{Q>0}
              < -\displaystyle\int \frac{\mathrm dP}{v}\big|_{a}.\)
  \end{itemize}
  \item Si \(Q<0\) (se enfría el gas):
  \begin{itemize}
    \item \(v\) es menor que en el adiabático.
    \item \(\Rightarrow -\displaystyle\int \frac{\mathrm dP}{v}\big|_{Q<0}
              > -\displaystyle\int \frac{\mathrm dP}{v}\big|_{a}.\)
  \end{itemize}
\end{itemize}

El logaritmo \(\ln(v_2/v_1)\) y el término de fricción no cambian mucho entre estos casos, de modo que:

\begin{equation}
\boxed{
\left(\frac{w}{A}\right)_{Q>0}
<
\left(\frac{w}{A}\right)_{\text{adiabático}}
<
\left(\frac{w}{A}\right)_{Q<0}
}
\end{equation}

Calentar el gas reduce el flujo respecto al adiabático; enfriarlo lo aumenta.

\bigskip
\hrule
\bigskip

\section*{4.\ Efecto del perfil de temperatura respecto al ``isotérmico a \(T_1\)''}

Ahora se toma como referencia un modelo isotérmico a \(T_1\), y se compara con
casos reales donde la temperatura a lo largo de la conducción es mayor o menor
que \(T_1\). De nuevo, gas ideal y mismo \((P_1,P_2)\):

\begin{itemize}
  \item Si en toda la tubería \(T > T_1 \Rightarrow v\) es mayor que en el isotérmico:
  \[
    -\int \frac{\mathrm dP}{v}\bigg|_{T>T_1}
    <
    -\int \frac{\mathrm dP}{v}\bigg|_{T=T_1}.
  \]
  \item Si en toda la tubería \(T < T_1 \Rightarrow v\) es menor:
  \[
    -\int \frac{\mathrm dP}{v}\bigg|_{T<T_1}
    >
    -\int \frac{\mathrm dP}{v}\bigg|_{T=T_1}.
  \]
\end{itemize}

Resultado:

\begin{equation}
\boxed{
\left(\frac{w}{A}\right)_{T>T_1}
<
\left(\frac{w}{A}\right)_{\text{isotérmico a }T_1}
<
\left(\frac{w}{A}\right)_{T<T_1}
}
\end{equation}

Si el gas está más caliente que \(T_1\) a lo largo de la línea, el flujo real
será menor que el del modelo isotérmico a \(T_1\); si está más frío, será mayor.

\bigskip
\hrule
\bigskip

\section*{5.\ Comparación de casos reales con modelos (mapa cualitativo)}

Las últimas ideas combinan todo lo anterior y ubican el flujo real
\(w_\text{real}\) entre los flujos de los modelos isotérmico (\(w_i\)) y
adiabático (\(w_a\)), según:

\begin{itemize}
  \item La relación entre \(T_1\) (temperatura de entrada) y la temperatura ambiente \(T_\text{amb}\).
  \item El signo de \(Q\) (calor hacia el gas o desde el gas).
  \item Si el perfil \(T(x)\) es globalmente mayor o menor que \(T_1\).
\end{itemize}

Los casos típicos:

\begin{enumerate}
  \item \textbf{Gas muy frío al inicio, ambiente caliente.}\\
        \(T_1 \ll T_\text{amb}\) y \(Q>0\), con \(T(x)>T_1\) a lo largo:
        \[
        w_\text{real} < w_i < w_a.
        \]
        El gas se calienta fuertemente: el flujo real queda por debajo del
        isotérmico, que a su vez está por debajo del adiabático.

  \item \textbf{Gas a temperatura similar a la ambiente, leve calentamiento/enfriamiento.}\\
        \(T_1 \approx T_\text{amb}\), \(Q>0\) moderado, con \(T(x)\) apenas por debajo de \(T_1\):
        \[
        w_i < w_\text{real} < w_a.
        \]
        El flujo real queda entre ambos modelos, con el isotérmico como cota
        inferior y el adiabático como cota superior.

  \item \textbf{Gas muy caliente, ambiente mucho más frío, enfriamiento intenso.}\\
        \(T_1 \gg T_\text{amb}\), \(Q<0\), con \(T(x)<T_1\) a lo largo:
        \[
        w_i < w_a < w_\text{real}.
        \]
        El gas se enfría mucho: el flujo real supera incluso al adiabático.

  \item \textbf{Casos intermedios.}\\
        Si \(T_1<T_\text{amb}\) y hay calentamiento pero no está claro el
        perfil respecto a \(T_1\): suele cumplirse \(w_\text{real} < w_a\).\\
        Si \(T_1>T_\text{amb}\) y el gas se mantiene más frío que \(T_1\):
        típicamente \(w_i < w_\text{real}\).
\end{enumerate}

De forma compacta, podemos resumir las desigualdades de referencia como:

\begin{align}
\left(\frac{w}{A}\right)_{\text{iso}}
&<
\left(\frac{w}{A}\right)_{\text{adiab}},
\\[4pt]
\left(\frac{w}{A}\right)_{Q>0}
&<
\left(\frac{w}{A}\right)_{\text{adiab}}
<
\left(\frac{w}{A}\right)_{Q<0},
\\[4pt]
\left(\frac{w}{A}\right)_{T>T_1}
&<
\left(\frac{w}{A}\right)_{\text{iso@}T_1}
<
\left(\frac{w}{A}\right)_{T<T_1}.
\end{align}

La posición de \(w_\text{real}\) respecto de \(w_i\) y \(w_a\) se deduce
combinando estas desigualdades con la información sobre \(T_1\), \(T_\text{amb}\)
y el signo de \(Q\).

\bigskip
\hrule
\bigskip

\section*{6.\ Criterio que voy a usar a partir de ahora}

Para todos los problemas de flujo compresible en conductos de sección
constante que resolvamos de aquí en más, adoptaré sistemáticamente estos
criterios:

\subsection*{6.1.\ Modelo adiabático}

\begin{itemize}
  \item Lo usaré cuando la conducción esté relativamente bien aislada o el
        tiempo de residencia sea corto.
  \item Lo utilizaré como \textbf{cota superior} de flujo másico (para un mismo
        \(\Delta P\)), salvo en casos de fuerte enfriamiento donde el real
        pueda superarlo (caso 3 de la sección anterior).
\end{itemize}

\subsection*{6.2.\ Modelo isotérmico (a \(T_1\) o a \(T_\text{amb}\))}

\begin{itemize}
  \item Apropiado cuando el tubo está bien acoplado térmicamente al ambiente y
        \(T_1 \simeq T_\text{amb}\).
  \item En general lo consideraré como \textbf{cota inferior} de flujo másico
        para un mismo \(\Delta P\).
\end{itemize}

\subsection*{6.3.\ Ubicación cualitativa del caso real}

\begin{itemize}
  \item Si \(Q>0\) (el gas se calienta globalmente), tenderé a ubicar
        \[
        w_\text{real} \le w_\text{adiab},
        \qquad
        \text{y frecuentemente } w_\text{real} \lesssim w_\text{iso}.
        \]
  \item Si \(Q<0\) (el gas se enfría),
        \[
        w_\text{real} \ge w_\text{adiab},
        \qquad
        \text{y típicamente } w_\text{real} \gtrsim w_\text{iso}.
        \]
  \item La relación de \(T_1\) con \(T_\text{amb}\) me dirá si
        \(w_\text{real}\) está por debajo de ambos modelos, entre ambos, o por
        encima de ambos, siguiendo los diagramas cualitativos (casos 1--3).
\end{itemize}

En resumen, usaré los modelos adiabático e isotérmico como modelos ideales
``faro'', y ubicaré el flujo real entre ellos usando estas desigualdades y
consideraciones térmicas cada vez que tengamos que elegir el modelo ideal o
justificar cuál usamos para acotar la realidad.

\end{document}
